\documentclass[11pt]{article}
\usepackage[margin=1in]{geometry}
\usepackage{amsmath,amssymb,amsthm}
\usepackage[hidelinks]{hyperref}

\newtheorem{theorem}{Theorem}
\newtheorem{remark}{Remark}

\newcommand{\R}{\mathbb{R}}
\newcommand{\C}{\mathbb{C}}
\newcommand{\Lp}{L^p}
\newcommand{\MW}{\widetilde M_W}
\newcommand{\Ap}{A_p}

\title{A Diagonal Exponent Improvement for the Matrix-Weighted Auxiliary Maximal Operator}
\author{Run fa\_banach\_001, agent\_lane\_02}
\date{2026-07-03}

\begin{document}
\maketitle

\section*{Source Question}

Kakaroumpas and Soler i Gibert prove in arXiv:2407.16776, Theorem 6, that for a
matrix $\Ap$ weight $W$ and $1<p,q<\infty$,
\[
 \left\|\left(\sum_{n=1}^{\infty}|\MW f_n|^q\right)^{1/q}\right\|_{\Lp(\R^N)}
 \le C [W]_{\Ap}^{\,\frac1p+\max\{\frac1{q-1},\frac1{p-1}\}}
 \left\|\left(\sum_{n=1}^{\infty}|W^{1/p} f_n|^q\right)^{1/q}\right\|_{\Lp(\R^N)} .
\]
Immediately before the theorem they note that they do not know whether this
exponent can be improved. The following observation gives such an improvement
on the diagonal $q=p$.

\section*{Proof Intuition}

The loss in Theorem 6 comes from a pointwise domination of $\MW f$ by a Hardy--Littlewood maximal
operator applied to the Christ--Goldberg maximal operator. On the diagonal $q=p$ no vector-valued
machinery is needed: the $L^p(\ell^p)$ norm is exactly the sum of the individual $L^p$ norms. Thus the
known one-function estimate for $\MW$ can be applied coordinate by coordinate.

\begin{theorem}
Let $1<p<\infty$, let $W$ be a $(d\times d)$ matrix $\Ap$ weight on $\R^N$, and let
$(f_n)$ be a finite or countable sequence of locally integrable $\C^d$-valued functions. Then
\[
 \left\|\left(\sum_{n=1}^{\infty}|\MW f_n|^p\right)^{1/p}\right\|_{\Lp(\R^N)}
 \le C_{N,d,p} [W]_{\Ap}^{\,1/(p-1)}
 \left\|\left(\sum_{n=1}^{\infty}|W^{1/p} f_n|^p\right)^{1/p}\right\|_{\Lp(\R^N)} .
\]
\end{theorem}

\begin{proof}
The source paper recalls, just before its vector-valued estimates, the one-function bound
\[
 \|\MW f\|_{\Lp(\R^N)}
 \le C_{N,d,p}[W]_{\Ap}^{1/(p-1)}
 \|\,|W^{1/p}f|\,\|_{\Lp(\R^N)} ,
\]
originally due to Isralowitz--Kwon--Pott for the auxiliary matrix-weighted maximal operator.

For a finite sequence $(f_n)_{n=1}^m$, raise the desired estimate to the $p$th power and use the
one-function estimate termwise:
\[
\begin{aligned}
 \left\|\left(\sum_{n=1}^{m}|\MW f_n|^p\right)^{1/p}\right\|_{\Lp}^p
 &= \sum_{n=1}^{m} \|\MW f_n\|_{\Lp}^p \\
 &\le C_{N,d,p}^p [W]_{\Ap}^{p/(p-1)}
       \sum_{n=1}^{m}\|\,|W^{1/p}f_n|\,\|_{\Lp}^p \\
 &= C_{N,d,p}^p [W]_{\Ap}^{p/(p-1)}
 \left\|\left(\sum_{n=1}^{m}|W^{1/p}f_n|^p\right)^{1/p}\right\|_{\Lp}^p .
\end{aligned}
\]
Taking $p$th roots proves the finite-sequence case. The countable case follows by monotone
convergence applied to the partial sums on both sides.
\end{proof}

\begin{remark}
For $q=p$, Theorem 6 in the source gives the larger exponent
\[
 \frac1p+\max\left\{\frac1{p-1},\frac1{p-1}\right\}
 = \frac1p+\frac1{p-1}.
\]
The theorem above removes the additional $1/p$ loss for all matrix $\Ap$ weights. It does not settle the
off-diagonal vector-valued estimates $q\ne p$.
\end{remark}

\section*{References}

\begin{itemize}
\item S. Kakaroumpas and O. Soler i Gibert, \emph{Vector valued estimates for matrix weighted maximal operators and product BMO}, arXiv:2407.16776.
\item I. Isralowitz, K. Kwon, and S. Pott, one-function bound for the Christ--Goldberg auxiliary maximal operator, as cited in arXiv:2407.16776.
\end{itemize}

\end{document}
